%% kapa: exam=1 round=28 year=2017 period=1 form=A subject=real_estate subject_ko=부동산학원론 seq=36 number=116
%% source: Q-Net 감정평가사 기출문제 (cst00309) artlSeq=5199603 file=extracted/1008/5199603/2017년 제28회 감정평가사 1차시험 문제지/1교시 A형.pdf
%% conversion: kordoc 4.24.51 -> tools/gen1.py (auto); review=proofread-sonnet; flags=none
\begin{question}{116}
다음은 부동산개발과정에 내재하는 위험에 관한 설명이다. ( )에 들어갈 내용으로 옳게 연결된 것은?
\begin{framedtext}
○ ( ㄱ )은 정부의 정책이나 용도지역제와 같은 토지이용규제의 변화로 인해 발생하기도 한다. \par
○ ( ㄴ )은 개발된 부동산이 분양이나 임대가 되지 않거나, 계획했던 가격 이하나 임대료 이하로 매각되거나 임대되는 경우를 말한다. \par
○ ( ㄷ )은 인플레이션이 심할수록, 개발기간이 연장될수록 더 커진다.
\end{framedtext}
\begin{choices}
  \item ㄱ: 법률적 위험, ㄴ: 시장위험, ㄷ: 비용위험
  \item ㄱ: 법률적 위험, ㄴ: 관리위험, ㄷ: 시장위험
  \item ㄱ: 사업위험, ㄴ: 계획위험, ㄷ: 비용위험
  \item ㄱ: 계획위험, ㄴ: 시장위험, ㄷ: 비용위험
  \item ㄱ: 시장위험, ㄴ: 계획위험, ㄷ: 사업위험
\end{choices}
\end{question}
\answer{1}
